\documentclass[10pt,letterpaper]{article}

% Packages
\usepackage[top=0.55in,bottom=0.55in,left=0.65in,right=0.65in]{geometry}
\usepackage{amsmath,amssymb}
\usepackage{microtype}
\usepackage{booktabs}
\usepackage{array}
\usepackage{xcolor}
\usepackage{hyperref}
\usepackage{enumitem}
\usepackage{titlesec}
\usepackage{tcolorbox}

% Color Palette
\definecolor{primary}{RGB}{15,44,89}
\definecolor{accent}{RGB}{20,110,180}
\definecolor{lightblue}{RGB}{239,246,252}
\definecolor{darktext}{RGB}{33,37,41}
\definecolor{muted}{RGB}{90,100,110}

\color{darktext}

% Hyperref
\hypersetup{
    colorlinks=true,
    linkcolor=accent,
    citecolor=accent,
    urlcolor=accent
}

% Section Formatting
\titleformat{\section}
    {\color{primary}\large\bfseries}
    {}{0em}
    {}
    [\vspace{-3pt}\color{accent}\rule{\textwidth}{0.7pt}]

\titlespacing*{\section}{0pt}{8pt}{4pt}

\titleformat{\subsection}
    {\color{primary}\normalsize\bfseries}
    {}{0em}{}

\titlespacing*{\subsection}{0pt}{4pt}{2pt}

% Paragraph / List Formatting
\setlength{\parindent}{0pt}
\setlength{\parskip}{3.5pt}

\setlist[itemize]{
    itemsep=1.5pt,
    topsep=2pt,
    parsep=0pt,
    leftmargin=1.3em
}

% Key Finding Box
\newtcolorbox{keybox}{
    colback=lightblue,
    colframe=accent,
    boxrule=0.5pt,
    arc=2pt,
    left=6pt,
    right=6pt,
    top=4pt,
    bottom=4pt
}

% Page Style
\pagestyle{empty}

% Document
\begin{document}

% TITLE
\begin{center}

{\LARGE\bfseries\color{primary}
Algorithmic Calibration of a JITAI Decision Engine}

\vspace{2pt}

{\large\bfseries\color{accent}
A Sensitivity Analysis Approach for Adaptive Micro-Randomized Trials}

\vspace{3pt}

{\normalsize\bfseries\color{primary}
Tien Le}

\vspace{2pt}

{\small
\textcolor{muted}{
Simulation-Based Calibration of Intervention Thresholds, Participant Burden, and EMA Compliance
}}

\end{center}

\vspace{-4pt}

% PURPOSE
\section{Purpose}

Just-In-Time Adaptive Interventions (JITAIs) aim to deliver timely behavioral nudges during acute moments of psychological vulnerability and emotional volatility. However, calibrating trigger thresholds presents a fundamental trade-off: over-sensitive thresholds induce notification fatigue, while overly conservative cutoffs miss critical therapeutic windows. This research introduces a simulation-based sensitivity framework to calibrate a randomized JITAI decision engine prior to clinical deployment.

The framework evaluates the relationship between intervention dosage, Ecological Momentary Assessment (EMA) compliance, and statistical efficiency. Our simulation identifies a \textbf{70\% compliance floor} below which parameter recovery degrades meaningfully, confirming a deliberate margin for the trial's pre-registered \textbf{75\% feasibility benchmark}, while calibrating an \textbf{80th-percentile within-person threshold} that yields an expected dosage of \textbf{4 decision points per week}.


{\footnotesize
\textbf{\color{primary}Keywords:}
Just-In-Time Adaptive Interventions (JITAI);
Micro-Randomized Trials (MRT);
Affective Computing;
Algorithmic Sensitivity Analysis;
Behavioral Informatics
}

% APPROACH
\section{Approach}

\subsection{Decision Engine}

The decision engine detects acute emotional volatility using the \textbf{Mean Squared Successive Difference (MSSD)}, computed over consecutive EMA observations (e.g., mood, stress, energy) within a rolling observation window:

\begin{equation}
\text{MSSD}
=
\frac{1}{n-1}
\sum_{i=1}^{n-1}
(x_{i+1}-x_i)^2
\end{equation}

Observed MSSD is evaluated against an empirical percentile threshold derived from the participant's historical data. The primary decision rule uses an \textbf{80th-percentile threshold}, a 3-observation window, a 60-minute cooldown, and a safety cap of four prompts daily. Observations exceeding this boundary trigger an eligible intervention decision point for micro-randomization within a Micro-Randomized Trial (MRT).

\subsection{Threshold Sensitivity Analysis}

To characterize how design choices affect intervention dosage, we evaluated a multi-parameter grid across MSSD percentiles, burn-in periods, and compliance rates.

\begin{center}
\footnotesize
\renewcommand{\arraystretch}{1.1}
\begin{tabular}{@{}lll@{}}
\toprule
\textbf{Parameter} & \textbf{Evaluated Grid Values} & \textbf{Design Objective} \\
\midrule
MSSD threshold ($\Theta$)
& 70th, 75th, 80th, 85th, 90th percentile
& Regulate prompt frequency \\

Burn-in period ($B$)
& 0, 1, 2, 3, 5, 7 days
& Evaluate baseline warm-up effects \\

EMA response rate ($R$)
& 70\%, 80\%, 90\%
& Characterize missing-data sensitivity \\
\bottomrule
\end{tabular}
\end{center}

To rigorously test the engine, the simulation framework modeled cohorts of 150 participants over 28 days, assuming \textbf{5 EMA opportunities per day as a fixed simulation parameter} (evaluating the 5--6 daily prompts scheduled in the clinical protocol). Across 10 random seeds, each threshold--burn-in combination was tested over \textbf{1,500 participant-level realizations}. The 80th-percentile threshold consistently produced a target dosage of \textbf{4-5 decision points per participant per week} following the burn-in period.

\subsection{Parameter Recovery \& Compliance Sensitivity}

A second sensitivity sweep evaluated latent parameter recovery across response rates from $70\%$ to $90\%$ (10\% increments; $N=1,500$ per level). Recovery of process volatility ($\sigma$) remained robust across missingness conditions. In contrast, recovering the temporal autocorrelation ($\rho$) proved highly sensitive to missing data: parameter recovery degraded sharply when compliance dropped below approximately \textbf{70\%}, but stabilized reliably above it, establishing \textbf{70\%} as the mathematical breakdown floor for temporal recovery. Consequently, both figures serve distinct roles: \textbf{75\%} represents the trial's pre-registered operational feasibility benchmark, which deliberately incorporates a \textbf{5\%} safety buffer above the \textbf{70\%} algorithmic floor to protect against inferential collapse. This demonstrates that participant compliance is substantially more critical for capturing temporal dependence than static volatility.

% IMPACT
\section{Impact \& Significance}

This framework replaces heuristic thresholding with an empirical, reproducible approach to participant burden and power calibration. By confirming that the 80th-percentile threshold stabilizes dosage at \textbf{4 decision points per week}, the analysis ensures sufficient longitudinal power to detect time-varying treatment effects in MRTs without inducing participant habituation. 

Furthermore, separating compliance effects on static variance ($\sigma$) versus temporal dependence ($\rho$) provides a generalizable methodological guideline for mobile health systems that pair active self-reports with continuous physiological telemetry (e.g. Heart Rate Variability).

\end{document}